\section{Численные результаты}
\label{sec:Chapter6} \index{Chapter6}

Вот результаты для $\mathsf{LAC}$ и $\mathsf{NewHope}$. Параметры $[n, q, k]$ были взяты $[1024, 251, 1]$ и $[1024, 12289, 8]$ соответственно. Целевой DFR положим $10^{-6}$.

Для $\mathsf{LAC}$ имеет смысл рассматривать размер алфавита $Q$ от $2$ до $6$ символов включительно. Посчитав распределение в $\Z_q$ шума на компьютере, бинпоиском подбираем необходимое кол-во исправляемых ошибок.

\begin{table}[h]
	\centering
	\begin{minipage}{0.48\textwidth}
		\centering
		\begin{tabular}{|c|c|c|}
			\hline Q & t & err \\ \hline
			2 & 21 & 0.0058 \\ \hline
			3 & 109 & 0.0668 \\ \hline
			4 & 235 & 0.1709 \\ \hline
			5 & 345 & 0.2695 \\ \hline
			6 & 448 & 0.3651 \\ \hline
		\end{tabular}
		\caption*{Hamming}
	\end{minipage}
	\begin{minipage}{0.48\textwidth}
		\centering
		\begin{tabular}{|c|c|c|c|}
			\hline Q & t & err1 & err2 \\ \hline
			2 & 21 & 0.0058 & 0 \\ \hline
			3 & 109 & 0.0668 & 0 \\ \hline
			4 & 235 & 0.1709 & 3.64e-05 \\ \hline
			5 & 346 & 0.2686 & 0.0009 \\ \hline
			6 & 456 & 0.3594 & 0.0058 \\ \hline
		\end{tabular}
		\caption*{Li}
	\end{minipage}
\end{table}

Последний столбец отвечает за вероятность ошибки в каждом символе (если рассматривать их независимо). В метрике Ли это вероятности отличия от исходного символа на расстояние $1$ и $2$ соответственно.

Для $\mathsf{NewHope}$ делаем то же самое, только рассматривая другие $Q$.

\begin{table}[h]
	\centering
	\begin{tabular}{|c|c|c|}
		\hline Q & t & err \\ \hline
		2 & 0 & $\approx 0$ \\ \hline
		3 & 0 & $\approx 0$ \\ \hline
		4 & 0 & $\approx 0$ \\ \hline
		5 & 0 & 1.52e-11 \\ \hline
		6 & 1 & 1.71e-08 \\ \hline
		7 & 1 & 1.32e-06 \\ \hline
		8 & 3 & 2.28e-05 \\ \hline
		9 & 5 & 1.66e-04 \\ \hline
	\end{tabular}	
	\caption*{Hamming, Li}
\end{table}

Для маленьких размеров алфавита вероятность двойной ошибки пренебрижимо мала, так что результаты одинаковы.

\begin{table}[h]
	\begin{minipage}{0.48\textwidth}
		\centering
		\begin{tabular}{|c|c|c|}
			\hline Q & t & err \\ \hline
			10 & 7 & 0.0007 \\ \hline
			12 & 18 & 0.0047 \\ \hline
			14 & 38 & 0.0155 \\ \hline
			16 & 65 & 0.0339 \\ \hline
			18 & 100 & 0.0596 \\ \hline
			20 & 138 & 0.0899 \\ \hline
			23 & 199 & 0.1402 \\ \hline
			26 & 259 & 0.1923 \\ \hline
			30 & 333 & 0.2585 \\ \hline
			34 & 398 & 0.3186 \\ \hline
			38 & 455 & 0.3722 \\ \hline
		\end{tabular}
		\caption*{Hamming}
	\end{minipage}
	\hfill
	\begin{minipage}{0.48\textwidth}
		\centering
		\begin{tabular}{|c|c|c|c|}
			\hline Q & t & err1 & err2 \\ \hline
			10 & 7 & 0.0007 & $\approx 0$ \\ \hline
			12 & 18 & 0.0047 & $\approx 0$ \\ \hline
			14 & 38 & 0.0155 & $\approx 0$ \\ \hline
			16 & 65 & 0.0339 & 2.32e-10 \\ \hline
			18 & 100 & 0.0596 & 1.74e-08 \\ \hline
			20 & 138 & 0.0899 & 3.87e-07 \\ \hline
			23 & 199 & 0.1402 & 1.00e-05 \\ \hline
			26 & 259 & 0.1922 & 0.0001 \\ \hline
			30 & 334 & 0.2578 & 0.0007 \\ \hline
			34 & 402 & 0.3158 & 0.0028 \\ \hline
			38 & 466 & 0.3647 & 0.0075 \\ \hline
		\end{tabular}
		\caption*{Li}
	\end{minipage}
\end{table}

Тем не менее, при больших значениях $Q$ разница проявляется. Дальше по известным $n$ и $d$ можно посчитать границы на скорость $R$.

\begin{table}[h]
	\centering
	\begin{minipage}{0.48\textwidth}
		\centering
		\begin{tabular}{|c|c|c|}
			\hline Q & $R_{GV}$ & $R_H$ \\ \hline
			2 & 0.7569 & 0.8591 \\ \hline
			3 & 0.6298 & 0.9939 \\ \hline
			4 & 0.2821 & 0.8639 \\ \hline
			5 & 0.0674 & 0.7312 \\ \hline
			6 & 0.0010 & 0.5854 \\ \hline
		\end{tabular}
		\caption*{LAC, Hamming}
	\end{minipage}
	\begin{minipage}{0.48\textwidth}
		\centering
		\begin{tabular}{|c|c|c|}
			\hline Q & $R_{GV}$ & $R_H$ \\ \hline
			2 & 0.7569 & 0.8591 \\ \hline
			3 & 0.6298 & 0.9939 \\ \hline
			4 & 0.4507 & 0.9767 \\ \hline
			5 & 0.3495 & 0.9804 \\ \hline
			6 & 0.3072 & 0.9879 \\ \hline
		\end{tabular}
		\caption*{LAC, Li}
	\end{minipage}
\end{table}

Видно, что при маленьких $Q$ и $t$ разница между метриками Хэмминга и Ли небольшая или вообще отсутствует, при увеличении же этих параметров Хэмминг начинает резко убывать, а Ли почти не уменьшается (за счёт гораздо меньшего объёма шаров, возникающих вокруг кодовых слов).

\newpage
\begin{table}[h]
	\centering
	\begin{minipage}{0.48\textwidth}
		\centering
		\begin{tabular}{|c|c|c|}
			\hline Q & $R_{GV}$ & $R_H$ \\ \hline
			2 & 1.0000 & 1.0000 \\ \hline
			3 & 1.5850 & 1.5850 \\ \hline
			4 & 2.0000 & 2.0000 \\ \hline
			5 & 2.3219 & 2.3219 \\ \hline
			6 & 2.5619 & 2.5729 \\ \hline
			7 & 2.7838 & 2.7951 \\ \hline
			8 & 2.9342 & 2.9650 \\ \hline
			9 & 3.0643 & 3.1132 \\ \hline
			10 & 3.1775 & 3.2439 \\ \hline
			12 & 3.2475 & 3.3999 \\ \hline
			14 & 3.1555 & 3.4450 \\ \hline
			16 & 2.9596 & 3.4151 \\ \hline
			18 & 2.6640 & 3.3137 \\ \hline
			20 & 2.3412 & 3.1837 \\ \hline
			23 & 1.8315 & 2.9513 \\ \hline
			26 & 1.3565 & 2.7150 \\ \hline
			30 & 0.8186 & 2.4222 \\ \hline
			34 & 0.4060 & 2.1680 \\ \hline
			38 & 0.1187 & 1.9473 \\ \hline
		\end{tabular}
		\caption*{NewHope, Hamming}
	\end{minipage}
	\begin{minipage}{0.48\textwidth}
		\centering
		\begin{tabular}{|c|c|c|}
			\hline Q & $R_{GV}$ & $R_H$ \\ \hline
			2 & 1.0000 & 1.0000 \\ \hline
			3 & 1.5850 & 1.5850 \\ \hline
			4 & 2.0000 & 2.0000 \\ \hline
			5 & 2.3219 & 2.3219 \\ \hline
			6 & 2.5645 & 2.5742 \\ \hline
			7 & 2.7868 & 2.7966 \\ \hline
			8 & 2.9448 & 2.9703 \\ \hline
			9 & 3.0838 & 3.1230 \\ \hline
			10 & 3.2070 & 3.2587 \\ \hline
			12 & 3.3331 & 3.4429 \\ \hline
			14 & 3.3518 & 3.5442 \\ \hline
			16 & 3.3163 & 3.5966 \\ \hline
			18 & 3.2367 & 3.6080 \\ \hline
			20 & 3.1565 & 3.6074 \\ \hline
			23 & 3.0427 & 3.5936 \\ \hline
			26 & 2.9581 & 3.5841 \\ \hline
			30 & 2.8863 & 3.5841 \\ \hline
			34 & 2.8495 & 3.5972 \\ \hline
			38 & 2.8290 & 3.6139 \\ \hline
		\end{tabular}
		\caption*{NewHope, Li}
	\end{minipage}
\end{table}

Также занимательно, что вероятность возникновения отличия на $\geq\!2$ в метрике Ли перестаёт быть пренебрежительно мала лишь к моменту, когда вероятность отличия на $1$ становится огромна (порядка $\frac14$ и выше).

Таким образом, пока кол-во исправляемых ошибок остаётся не слишком большим, можно считать их все переходами из символа в соседний.
